% !TEX program = xelatex
% DeepMIR HW2 report slides. Build: latexmk -xelatex main.tex   (16:9, theme = Moloch, the maintained Metropolis)
\documentclass[aspectratio=169,11pt]{beamer}
\usetheme[progressbar=frametitle]{moloch}        % swap for metropolis / focus / Arguelles to change the look

% ---------- fonts: Fira Sans (ships with TeX Live) + PingFang TC for Chinese ----------
\usepackage{fontspec}
\setsansfont{FiraSans}[Extension=.otf, UprightFont=*-Regular, BoldFont=*-SemiBold, ItalicFont=*-Italic, BoldItalicFont=*-SemiBoldItalic]
\setmonofont{FiraMono}[Extension=.otf, UprightFont=*-Regular, BoldFont=*-Bold, Scale=0.9]
\usepackage{xeCJK}
\setCJKmainfont{PingFangTC-Regular}[BoldFont=PingFangTC-Semibold]
\setCJKsansfont{PingFangTC-Regular}[BoldFont=PingFangTC-Semibold]
\setCJKmonofont{PingFangTC-Regular}

% ---------- packages ----------
\usepackage{booktabs, graphicx, array, multirow}
\usepackage{appendixnumberbeamer}
\graphicspath{{figs/}{../results/}}
\hypersetup{colorlinks=true, urlcolor=orange!80!black, linkcolor=.}

% ---------- helpers ----------
\newcommand{\best}[1]{\textbf{\alert{#1}}}                       % highlight the best number in a table
\newcommand{\fig}[2][0.78\textheight]{\includegraphics[height=#1,keepaspectratio]{#2}}
\usetikzlibrary{positioning, arrows.meta, calc}
\tikzset{
  blk/.style={draw=black!35, rounded corners=2pt, align=center, font=\scriptsize, inner sep=3.5pt, minimum height=1.0cm},
  io/.style={blk, fill=white},
  op/.style={blk, fill=black!5},
  train/.style={blk, fill=orange!22, draw=orange!75!black, thick},
  arr/.style={-{Stealth[length=4.5pt]}, semithick, black!65},
}
\newenvironment{tight}{\begin{itemize}\setlength{\itemsep}{1.5pt}\setlength{\parskip}{0pt}\setlength{\topsep}{0pt}}{\end{itemize}}
\newcommand{\src}[1]{\vfill{\scriptsize\color{gray}#1}}          % small source / setting note at the bottom of a slide
\newcommand{\todo}[1]{{\color{red}\textbf{TODO:} #1}}

\title{Symbolic Music Generation with REMI Transformers}
\subtitle{DeepMIR 2026 Fall --- Homework 2}
\author{Chi-Jen Peng \quad (R14725022)}
\institute{Code, README, checkpoints, generated samples: \todo{cloud link}\\[0.3em]
Repository: \url{https://github.com/chijenpeng/26-fall-DeepMIR-HW2}}
\date{October 29, 2026}

\begin{document}
\maketitle

% =====================================================================
\section{Experimental Setting}
\begin{frame}{Experimental setting}
  \begin{columns}[T]
    \begin{column}{0.5\textwidth}
      \textbf{Tasks}
      \begin{tight}
        \item Basic A: unconditional 32-bar generation, 3 inference settings $\times$ 20 pieces
        \item Basic B: 24-bar continuation of 3 TA prompts (8 bars), 3 settings
        \item Bonus (Option B): piano accompaniment for 10 melody prompts (16 bars)
        \item Data: all 896 POP909 training songs, no other data; models trained from scratch
      \end{tight}
    \end{column}
    \begin{column}{0.5\textwidth}
      \textbf{Hardware / software}\\[0.3em]
      {\small
      \begin{tabular}{@{}ll@{}}
        \toprule
        GPU & 2 $\times$ RTX 4070 SUPER (12 GB),\\ & 1 $\times$ RTX A4500 (20 GB) \\
        OS & Ubuntu 26.04 \\
        Software & PyTorch 2.8, Transformers 4.57.1,\\ & pretty\_midi 0.2.11, FluidSynth \\
        \bottomrule
      \end{tabular}}
    \end{column}
  \end{columns}
\end{frame}

% =====================================================================
\section{Method}
\begin{frame}[t]{Token representation: REMI on the annotated beat grid}
  % TODO: example token sequence figure / table
  \begin{tight}
    \item Bars start at the downbeats in \texttt{beat\_midi.txt}; each bar: \texttt{Bar Beat\_n Tempo\_t}
    \item Events: \texttt{Pos\_p} (16th note, 4 per beat) \texttt{[Chord\_root:quality]} \texttt{Pitch Vel Dur}
    \item Vocabulary: 430 tokens; 32 bars: median 2{,}431 tokens, p99 3{,}900
    \item \todo{chords, velocity / duration bins, Basic vs.\ Bonus tracks}
  \end{tight}
\end{frame}

\begin{frame}[t]{Models and training}
  \todo{GPT-2 (Basic) and T5 (Bonus) architecture, hyper-parameters, augmentation}
\end{frame}

\begin{frame}[t]{Training length: held-out loss}
  \todo{held-out loss vs.\ step (13 prompt songs), choice of checkpoints}
\end{frame}

% =====================================================================
\section{Basic Task A: Unconditional Generation}
\begin{frame}[t]{Basic A: inference settings and objective metrics}
  \centering\small
  \begin{tabular}{@{}lllcccccc@{}}
    \toprule
    Model & Repr. & w/ Chord & Top-K / Top-P & Temp. & H1 & H4 & GS \\
    \midrule
    GPT-2 & REMI & O & K = 20 & 0.9 & -- & -- & -- \\
    GPT-2 & REMI & O & P = 0.9 & 1.0 & -- & -- & -- \\
    GPT-2 & REMI & O & K = 50 & 1.2 & -- & -- & -- \\
    GPT-2 & REMI & X & K = 20 & 0.9 & -- & -- & -- \\
    \midrule
    \multicolumn{5}{c}{Real data (896 songs)} & 2.123 & 2.626 & 0.877 \\
    \bottomrule
  \end{tabular}
  \src{20 pieces $\times$ 32 bars per row; metrics from the TA \texttt{eval\_basic.py} (MusDr).}
\end{frame}

% =====================================================================
\section{Basic Task B: Prompt Continuation}
\begin{frame}[t]{Basic B: continuation of the three prompts}
  \todo{settings, best setting per prompt (listening)}
\end{frame}

% =====================================================================
\section{Bonus: Accompaniment Generation}
\begin{frame}[t]{Bonus: melody $\rightarrow$ piano with T5}
  \todo{model, inference, CA / OA table}
\end{frame}

% =====================================================================
\section{Discussion}
\begin{frame}[t]{Discussion and limitations}
  \todo{findings}
\end{frame}

% =====================================================================
\section{References}
\begin{frame}[t]{References}
  \footnotesize
  \begin{tight}\footnotesize
    \item \textbf{POP909}: Z.~Wang et al., ``POP909: A Pop-song Dataset for Music Arrangement Generation,'' ISMIR 2020
    \item \textbf{REMI}: Y.-S.~Huang, Y.-H.~Yang, ``Pop Music Transformer,'' ACM MM 2020
    \item \textbf{MusDr} metrics: S.-L.~Wu, Y.-H.~Yang, ``The Jazz Transformer on the Front Line,'' ISMIR 2020
    \item \textbf{Bonus metrics}: Y.~Ren et al., ``PopMAG: Pop Music Accompaniment Generation,'' ACM MM 2020
    \item \textbf{GPT-2}: A.~Radford et al., 2019. \quad \textbf{T5}: C.~Raffel et al., JMLR 2020
    \item Hugging Face Transformers: T.~Wolf et al., EMNLP 2020. \quad PyTorch: A.~Paszke et al., NeurIPS 2019
    \item TA scripts (\texttt{tokenization.py}, \texttt{eval\_basic.py}, \texttt{eval\_bonus.py}); TimGM6mb SoundFont (GPL-2)
  \end{tight}
  \vspace{0.5em}
  \textbf{Tool}: Claude Code (Anthropic), OpenAI Codex.
\end{frame}

\end{document}
